The replacement of technical specification 0126 with specification 1580 changed Enel Distribuição Ceará's requirements for connecting buildings with multiple consumer units. This study compares the requirements applicable in Ceará and assesses their implications for the approved electrical design of the Palatium residential building in Sobral, Brazil. The two specifications were examined alongside the design calculations and drawings. Of the 82 requirements compared, 39 were selected for the case study, yielding 38 impact records. Of these, 23 indicated potentially necessary changes, seven indicated compatibility, two required assessment by the utility, and six could not be resolved from the available documents. Under the scenario examined, relocating the main protection assembly affects the location of fire protection system metering and surge protection. The existing electrical branch supplying the fire protection system can be retained. The reported demand of 291.33~kVA was used as the reference because the documents neither quantify the contribution of electric vehicle charging to coincident demand nor explain the residential demand multiplier. These findings identify the design decisions requiring review and the information needed to resolve the outstanding cases.

\keywords{electric power distribution; multi-unit service connections; electrical demand; electricity metering; electrical design.}
